\chapter{两种基于动态延迟主元的波前矩阵分解Kernel}
\label{chap:chap3}
在超节点方法中，需要把单个超节点的主元块作为稠密矩阵进行 $LU$ 或 $L(D)L^T$ 分解。
按照列主元方法，算法逐列选择主元并进行更新。如果某一列待分解的部分全为零，就无法选出列主元。
LAPACK 的 \code{dgetrf} 函数在大量的求解器中担负这个责任,
但当其遇到上述问题的时候则会跳过该零列，继续分解后续列并在最后报告错误；
这种处理不能满足我们对不正定但是满秩可分解矩阵的分解流程的需要。

一个子主元块无法分解，并不代表整个矩阵无法分解。
当前节点不能消去的部分可以传递给父节点，再在父节点中尝试分解。

需要处理的波前矩阵由四部分组成：

\begin{figure}[htbp]
\centering
\begin{tikzpicture}[x=0.82cm,y=0.82cm,font=\small]
  \fill[mumps fully summed] (0,3.0) rectangle (4.4,7.2);
  \fill[mumps factor] (4.4,3.0) rectangle (7.2,7.2);
  \fill[mumps factor] (0,0) rectangle (4.4,3.0);
  \fill[mumps contribution] (4.4,0) rectangle (7.2,3.0);
  \draw[thick] (0,0) rectangle (7.2,7.2);
  \draw[thick] (4.4,0)--(4.4,7.2);
  \draw[thick] (0,3.0)--(7.2,3.0);
  \node at (2.2,5.1) {Fully Summed};
  \node at (5.8,5.1) {U Factor};
  \node at (2.2,1.5) {L Factor};
  \node at (5.8,1.5) {CB Block};
\end{tikzpicture}
\caption{非对称超节点对应的波前矩阵}
\label{fig:appendix-a-unsym-front}
\end{figure}
\begin{itemize}
  \item Fully Summed 是已经装配完成、可以进行分解的矩阵，即超节点的对角块；
  \item L/U Factor 是需要根据 Fully Summed 的分解结果更新的块，最终作为因子保存；
  \item CB Block 是子节点传给父节点的贡献块，需要在更新后传递给父节点；
  \item 对角块之外的行列来自与对角块中的行列相耦合的变量。
\end{itemize}

我们针对上述问题提出了两种基于动态延迟主元的波前矩阵分解Kernel，
分别用于处理对称和非对称稠密矩阵。

\section{非对称稠密矩阵的延迟主元 LU 方法}
\subsection{分解 Fully Summed 块}

首先单独考虑 Fully Summed 块的分解。
已经分解更新的部分包含 $L_{EE}$、$U_{EE}$、$L_{DE}$ 和 $U_{ED}$，
接下来需要继续分解 Unfactored Block。一个朴素的想法是：
首先选择列主元（即该列待更新的行中最大元素），如果能够选到合适的列主元，
则进行行交换 $P$，并用该元素作为主元；
如果不能选到合适的列主元，则在整个活动子块中选择主元（即待更新的子块中最大的元素），
此时需要同时进行行、列交换，将活动子块的全局最大元素置换到当前主元位置，
根据该主元计算 $L_p$，再更新 $U_p$ 和右下角子块。

\begin{figure}[htbp]
\centering
\begin{tikzpicture}[x=0.82cm,y=0.82cm,font=\small]
  % 分解当前主元之前
  \begin{scope}
    \fill[mumps fully summed] (0,4) rectangle (2,6);
    \fill[mumps factor] (2,4) rectangle (6,6);
    \fill[mumps factor] (0,0) rectangle (2,4);
    \fill[mumps active] (2,0) rectangle (6,4);
    \draw[thick] (0,4) rectangle (2,6);
    \draw[thick] (0,4)--(6,4);
    \draw[thick] (2,0)--(2,6);
    \draw[thick] (0,0) rectangle (6,6);
    \draw[dashed] (0,6)--(2,4);
    \node at (1.42,5.30) {$U_{EE}$};
    \node at (0.66,4.66) {$L_{EE}$};
    \node at (4,5) {$U_{ED}$};
    \node at (1,2) {$L_{DE}$};
    \node[align=center] at (4,2) {Unfactored\\Block};
  \end{scope}

  % 接受一个主元之后
  \begin{scope}[xshift=7.0cm]
    \fill[mumps fully summed] (0,4) rectangle (2,6);
    \fill[mumps factor] (2,4) rectangle (6,6);
    \fill[mumps factor] (0,0) rectangle (2,4);
    \fill[mumps active] (2,0) rectangle (6,4);
    \fill[mumps delayed] (2,3.1) rectangle (2.9,4);
    \fill[mumps factor] (2.9,3.1) rectangle (6,4);
    \fill[mumps factor] (2,0) rectangle (2.9,3.1);
    \draw[thick] (0,4) rectangle (2,6);
    \draw[thick] (0,4)--(6,4);
    \draw[thick] (2,0)--(2,6);
    \draw[thick] (0,0) rectangle (6,6);
    \draw[dashed] (0,6)--(2,4);
    \node at (1.42,5.30) {$U_{EE}$};
    \node at (0.66,4.66) {$L_{EE}$};
    \node at (4,5) {$U_{ED}$};
    \node at (1,2) {$L_{DE}$};

    \draw (2.9,0)--(2.9,4);
    \draw (2,3.1)--(6,3.1);
    \node at (2.45,3.55) {$P$};
    \node at (4.45,3.55) {$U_P$};
    \node at (2.45,1.55) {$L_P$};
    \node at (4.45,1.55) {$\cdots$};
  \end{scope}
\end{tikzpicture}
\caption{非对称 Fully Summed 块的一步主元分解过程}
\label{fig:unsym-pivot-step}
\end{figure}
\subsubsection{主元接受阈值}
设第 $k$ 步尚未分解的子块为
$S^{(k)}=\left[s_{ij}^{(k)}\right]_{i,j=k}^{f}$，
其中 $f$ 为 Fully Summed 块的阶数，$\tau\in(0,1)$ 为主元接受阈值。
列主元只搜索活动子块第 $k$ 列，候选列主元所在行为：
$$p_k=\underset{k\le i\le f}{\operatorname{argmax}}
      \left|s_{ik}^{(k)}\right|,
\qquad
\alpha_k=\left|s_{p_k k}^{(k)}\right|.$$
\qquad 记当前子块的侯选主元为 $h_k$：
当
$h_k\ge\tau\alpha_k$
时，认为当前的侯选主元主元足够稳定，于是不用做交换就可以利用其作为主元进行分解。反之 $h_k<\tau\alpha_k$ 时，说明
需要将侯选列主元选作主元，即将当前行与第P行作交换。

但是还存在一种特殊情况：当前列为全零列，选出的侯选主元和侯选列主元都为零，就无法选出列主元。
此时需要在整个活动子块中搜索最大元素作为主元。
记活动子块中的最大元素为：
$$\beta_k=\left\|S^{(k)}\right\|_{\max}
=\max_{k\le i,j\le f}\left|s_{ij}^{(k)}\right|,\qquad
(p_k,q_k)=\underset{k\le i,j\le f}{\operatorname{argmax}}
           \left|s_{ij}^{(k)}\right|.
$$
分别用 $P_k$ 和 $Q_k$ 将第 $p_k$ 行、第 $q_k$ 列交换到第 $k$ 行、第 $k$ 列：
$$\widetilde S^{(k)}=P_kS^{(k)}Q_k,
\qquad
d_k=\widetilde s_{kk}^{(k)}=s_{p_kq_k}^{(k)},
\qquad
|d_k|=\beta_k.$$
因此，只使用列主元时只需要记录行置换，最终分解满足 $PA=LU$；如果使用过全主元，则必须记录行、列两种置换，最终满足
$PAQ=LU$。

\subsubsection{列延迟优化}
直接进行全主元搜索和逐项更新会消耗大量时间，这样的性能差距是我们无法接受的，所以我们提出了列延迟优化方法来取代全主元搜索。

全主元通过行列变换把全局最大元素移动到当前主元位置，
本质上是先把全局最大元素所在列交换到第 $k$ 列，再在换入列中进行一次列主元分解，
相当于提前对这一列进行了选主元分解的步骤。
于是我们可以反其道而行之，在原来的第 $k$ 列在活动部分全为零时，可以直接移到未分解区域尾部。
被移动到右侧的列在当前活动部分满足：
$$A(k:m-1,j)=0.$$
而后续 $LU$ 更新为：
$$A(k+1:m-1,j)
\leftarrow
A(k+1:m-1,j)-L(:,k)U(k,j).$$
由于
$U(k,j)=0$，
更新项也为零，该列在后续分解中仍保持为零。因此，全主元步骤可以替换为以下列延迟步骤：

\begin{enumerate}
  \item 若 $\alpha_k=0$，将该列移动到 Unfactored Block 的尾部，并记录此次列交换；
  \item 从右向左寻找第一列可用列（为防止已进入尾部的零列被多次交换，搜索范围必须排除已形成的延迟列），找到后立即停止，不再比较其他列；
  \item 在该列内部使用列主元法，并将该列与当前被延迟的列交换；
  \item 如果搜索到当前活动区的起点仍找不到可用列，说明整个 Unfactored Block 已经全零，判定为将剩余主元全部延迟。
\end{enumerate}
\subsubsection{right-looking 更新}
主元 $d_k$ 确定后，对主元行、主元列和右下角剩余子块执行 right-looking 更新：
$$(L_p)_i=\frac{\widetilde s_{ik}^{(k)}}{d_k},
\quad k<i\le f,
\qquad
(U_p)_j=\widetilde s_{kj}^{(k)},
\quad k<j\le f,$$
$$s_{ij}^{(k+1)}
=\widetilde s_{ij}^{(k)}-(L_p)_i(U_p)_j
=\widetilde s_{ij}^{(k)}
-\frac{\widetilde s_{ik}^{(k)}\widetilde s_{kj}^{(k)}}{d_k},
\qquad k<i,j\le f.$$
\qquad 若 $\beta_k=0$，则 $S^{(k)}$ 是全零块，列主元和全主元均不存在，在数值意义下不能继续分解，需要将未分解变量延迟到父节点。

完成上述步骤后，$P$、$U_p$ 和 $L_p$ 归入已分解部分，剩余部分继续作为 Unfactored Block 进行下一步迭代。
\subsubsection{自适应 KJI--SAXPY 分块优化}
\textbf{这里可以引用下那篇论文的KJI-SAXPY方法}

如果每次选择主元后立即更新所有剩余元素，计算会包含大量细粒度操作。
受到LAPACK 的 \code{dgetrf} 函数的起发，
我们选择借助 cuBLAS 三角求解和矩阵乘法对矩阵块进行集中更新，
即先分解一部分主元，再统一更新剩余部分。

示例选择的 panel 长度为 64。在 panel 内选择列主元时，
搜索范围包含当前64列的全部活动行的长方形矩阵。选择主元后，
只更新当前长方形 panel 中的数值；只要保证下一次选择主元之前当前列已经包含此前主元的更新即可。
完成一个 panel 后，再更新其余部分并开始新的 panel。

延迟主元会影响正在处理的 panel。为了保证数值分解正确，
算法在遇到延迟主元时立即停止当前 panel 的分解，
根据当前 panel 已经接受的主元把矩阵剩余部分完整更新一次，
然后从交换进来的非零列开始新的 panel。这样既使用 cuBLAS 加速了分解，
又保证主元搜索作用于最新的 Schur 补。

测试中，列延迟方法使整个计算过程的时间缩短约 $90\%$；
采用更细粒度的自适应 KJI--SAXPY 方法后，速度还能提升约一倍。整体优化后的算子在经过特别设计的算例上的运行时间已经略快于 LAPACK。

\subsection{更新波前矩阵}

Fully Summed 部分完成分解后，波前矩阵按已分解变量 $E$、延迟变量 $D$ 和普通更新变量 $U$ 分成以下部分：

\begin{figure}[htbp]
\centering
\begin{tikzpicture}[x=0.84cm,y=0.84cm,font=\small]
  \def\e{3.6}
  \def\d{4.9}
  \def\t{7.4}
  \fill[mumps fully summed] (0,3.7) rectangle (\e,\t);
  \fill[mumps delayed] (\e,3.7) rectangle (\d,\t);
  \fill[mumps factor] (\d,3.7) rectangle (\t,\t);
  \fill[mumps delayed] (0,2.3) rectangle (\t,3.7);
  \fill[mumps factor] (\d,2.3) rectangle (\t,3.7);
  \fill[mumps factor] (0,0) rectangle (\t,2.3);
  \draw[thick] (0,0) rectangle (\t,\t);
  \draw[thick] (\e,0)--(\e,\t);
  \draw[thick] (\d,0)--(\d,\t);
  \draw[thick] (0,2.3)--(\t,2.3);
  \draw[thick] (0,3.7)--(\t,3.7);
  \draw[dashed] (0,\t)--(\e,3.7);
  \node at (1.2,4.85) {$L_{EE}$};
  \node at (2.6,6.05) {$U_{EE}$};
  \node at (4.25,5.55) {$U_{ED}$};
  \node at (6.15,5.55) {$A_{EU}$};
  \node at (1.8,3.0) {$L_{DE}$};
  \node at (4.25,3.0) {$0$};
  \node at (6.15,3.0) {$A_{DU}$};
  \node at (1.8,1.15) {$A_{UE}$};
  \node at (4.25,1.15) {$A_{UD}$};
  \node at (6.15,1.15) {$A_{UU}$};
  \node[above] at (1.8,7.48) {$E$};
  \node[above] at (4.25,7.48) {$D$};
  \node[above] at (6.15,7.48) {$U$};
\end{tikzpicture}
\caption{Step 1 完成后波前矩阵的分区。}
\label{fig:appendix-a-unsym-step2}
\end{figure}
\textbf{数学推导？}

首先更新 $A_{EU}$ 和 $A_{UE}$ 为 $U_{EU}$ 和 $L_{UE}$，利用三角求解计算：
$$U_{EU}=L_{EE}^{-1}A_{EU}
\qquad
L_{UE}=A_{UE}U_{EE}^{-1}.$$
其次需要更新：$A_{DU}$、$A_{UD}$ 和 $A_{UU}$，分别计算：
$$A_{DU}\leftarrow A_{DU}-L_{DE}U_{EU}\qquad
A_{UD}\leftarrow A_{UD}-L_{UE}U_{ED}\qquad
A_{UU}\leftarrow A_{UU}-L_{UE}U_{EU}.
$$
\subsection{传递贡献块给父节点}

子节点的波前矩阵更新完成后，将右下的四个子矩阵组成扩展贡献块，并传递给父节点装配。
装配过程需根据贡献块的全局行、列编号执行 extend-add，
普通更新变量与父节点已有索引匹配后进行累加；
子节点未消去的延迟变量则需要动态扩充父节点的候选主元区域，
将其规划进父节点的 Fully Summed 块中。

\begin{figure}[htbp]
\centering
\begin{tikzpicture}[x=0.78cm,y=0.78cm,font=\small]
  % 父节点的整个波前矩阵
  \fill[mumps fully summed] (0,3) rectangle (7,10);
  \fill[mumps factor] (7,3) rectangle (10,10);
  \fill[mumps factor] (0,0) rectangle (7,3);
  \fill[mumps contribution] (7,0) rectangle (10,3);
  \draw[thick] (0,0) rectangle (10,10);

  % 已处理的超节点 S_1,\ldots,S_n
  \draw[thick,mumps fully summed] (0,9.3) rectangle (0.7,10);
  \node at (0.35,9.65) {$S_1$};
  \node at (1.10,8.95) {$\cdots$};
  \draw[thick,mumps fully summed] (1.5,8.05) rectangle (2.2,8.75);
  \node at (1.85,8.40) {$S_n$};

  % 超节点行、列边界的延伸线
  \draw[dashed] (0.7,9.3)--(10,9.3);
  \draw[dashed] (0.7,9.3)--(0.7,0);
  \draw[dashed] (1.5,8.05)--(1.5,0);
  \draw[dashed] (2.2,8.75)--(10,8.75);

  % 父波前右下区域及扩展贡献块边界
  \draw[thick] (2.2,8.05)--(10,8.05);
  \draw[thick] (2.2,8.05)--(2.2,0);
  \draw[mumps emphasis] (7,10)--(7,3);
  \draw[mumps emphasis] (0,3)--(7,3);
  \draw[thick] (7,3)--(7,0);
  \draw[thick] (7,3)--(10,3);

  % 分块名称
  \node[align=center] at (4.6,5.5) {Fully Summed\\[2pt]$E_p\cup D_c$};
  \node[align=center] at (8.5,5.5) {U Factor};
  \node[align=center] at (4.6,1.5) {L Factor};
  \node at (8.5,1.5) {CB Block};
\end{tikzpicture}
\caption{延迟主元贡献块向父节点传递}
\label{fig:appendix-a-extend-add}
\end{figure}

设子节点已消去的变量为 $E_c$，未能消去的延迟变量为 $D_c$，原来的普通更新变量为 $U_c$。
则子节点传出的扩展贡献块为：
$$C_c=
\begin{bmatrix}
S_{DD}&S_{DU}\\
S_{UD}&S_{UU}
\end{bmatrix},
\qquad
I_c^{\mathrm{out}}=D_c\cup U_c.$$

除了数值 $C_c$，还必须同时传递其全局行、列索引数组。对于非对称 LU，
行索引和列索引应分别保存：
$R_c^{\mathrm{out}}=D_c^r\cup U_c^r$、$
C_c^{\mathrm{out}}=D_c^c\cup U_c^c$。

设父节点原来的候选主元集为 $E_p$，待更新集为 $U_p$。
正确的符号分析通常保证了子节点的普通更新变量已经位于父波前中，即：
$$U_c\subseteq E_p\cup U_p.$$
子节点的普通更新变量到达父节点后，需要根据符号分析确定的消去归属重新分类：
\begin{itemize}
  \item $U_c\cap E_p$ 在子节点中只是更新变量，但到达父节点后已经 Fully Summed，
  可以在父节点中作为侯选部分参与主元选择；
  \item $U_c\cap U_p$ 仍然位于父节点的待更新区。
\end{itemize}
延迟变量 $D_c$ 是数值分解过程中动态产生的，父节点需要将其加入新的候选主元集：
$$E_p^{\mathrm{new}}=E_p\cup D_c,
\qquad
I_p^{\mathrm{new}}=E_p^{\mathrm{new}}\cup U_p.$$

得到了子节点传过来贡献块数据后，父节点为每个全局行、列编号建立到局部存储位置的映射：
$$\operatorname{map}_r(g)
=g\text{ 在父波前行索引中的局部位置}.$$
$$\operatorname{map}_c(g)
=g\text{ 在父波前列索引中的局部位置}.$$
\qquad 如果扩展后出现新的行或列，先扩展父波前的存储空间，并把新位置初始化为零。
然后再将所有子贡献块中全局编号一致的位置累加到父波前的新位置上：
$$F_p\!\left(
\operatorname{map}_r(R_c^{\mathrm{out}}(i)),
\operatorname{map}_c(C_c^{\mathrm{out}}(j))
\right)
\mathrel{+}=C_c(i,j).$$
\qquad 父矩阵中原来没有数值的位置按零处理，累加子节点贡献后成为新的填充。
如果多个子节点对同一全局位置都有贡献，必须把这些贡献全部累加，不能相互覆盖。

完成所有子贡献块的 extend-add 后，父节点按照以下规则处理：

\begin{itemize}
  \item $E_p\cup D_c$ 放入 Fully Summed 候选主元区，在父节点重新执行主元检验；
  \item 在父节点重复上述的选主元和更新步骤，成功消去的延迟变量进入 $L$、$U$ 因子；
  \item 仍无法通过主元检验的变量再次进入父节点的贡献块，继续向更高层节点传递。
\end{itemize}


\section{对称稠密矩阵的延迟 Bunch--Kaufman 方法}

如果整个波前矩阵是对称的，则 Fully Summed 块也是对称的。
此时采用 Bunch--Kaufman 方法进行 $LDL^T$ 分解：
$$P^TAP=LDL^T.$$
其中，$P$ 是同时置换行、列的置换矩阵，$L$ 是单位下三角矩阵，
$D$ 是由 $1\times1$ 和 $2\times2$ 主元块组成的块对角矩阵。
与非对称 LU 不同，对称分解不能只交换行而不交换列；交换第 $i$、$j$ 行时，
必须同时交换第 $i$、$j$ 列。

所以我们在对对称矩阵处理时主元块分解的方法和非对称矩阵不同，而剩余的更新步骤和非对称情况保持一致。
\subsection{Bunch--Kaufman 选主元}

设第 $k$ 步尚未分解的活动对称子块为
$S^{(k)}=\left[s_{ij}^{(k)}\right]_{i,j=k}^{f}$，
其中 $f$ 为 Fully Summed 块的阶数。
Bunch--Kaufman 方法在每一步选择一个 $1\times1$ 或 $2\times2$ 对角主元块，
并保持后续更新仍然对称。
定义 Bunch--Kaufman 阈值如下：
$$\gamma=\frac{1+\sqrt{17}}{8}\approx0.640388.$$
  \begin{figure}[htbp]
  \centering
  \begin{tikzpicture}[x=0.75cm,y=0.75cm,font=\small]
    % 已分解部分与当前活动对称子块
    \fill[mumps fully summed] (0,8) rectangle (2,10);
    \fill[mumps factor] (2,8) rectangle (10,10);
    \fill[mumps factor] (0,0) rectangle (2,8);
    \fill[mumps active] (2,0) rectangle (10,8);
    \fill[mumps delayed] (2,3.25) rectangle (10,4.25);
    \fill[mumps delayed] (2.75,0) rectangle (3.75,8);
    \draw[thick] (0,0) rectangle (10,10);
    \draw[thick] (2,0)--(2,10);
    \draw[thick] (0,8)--(10,8);
    \draw[dashed,gray] (0,10)--(2,8);
    \draw[dashed,gray] (2,8)--(10,0);

    \node at (1.42,9.42) {$U$};
    \node at (0.58,8.58) {$L$};
    \node at (6,9) {$L^T$ Factor};
    \node[rotate=90] at (1,4) {$L$ Factor};

    % k、p、j 所在行列的辅助线
    \draw[densely dashed,gray] (3.25,0)--(3.25,8);
    \draw[densely dashed,gray] (2,3.75)--(10,3.75);
    \draw[densely dashed,gray] (6.25,0)--(6.25,8);
    \draw[densely dashed,gray] (8.35,0)--(8.35,8);

    % 行列索引
    \node[above] at (3.25,8) {$k$};
    \node[above] at (6.25,8) {$p$};
    \node[above] at (8.35,8) {$j$};
    \node[left] at (2,3.75) {$p$};

    % Bunch--Kaufman 判定涉及的四个元素
    \filldraw[fill=MumpsBlue,draw=MumpsBlue] (3.25,6.75) circle (2.5pt);
    \node[anchor=south west] at (3.38,6.88)
      {$s_{kk}^{(k)}$};

    \filldraw[fill=MumpsRed,draw=MumpsRed] (3.25,3.75) circle (2.5pt);
    \node[anchor=north east] at (3.12,3.62)
      {$s_{pk}^{(k)}$};

    \filldraw[fill=MumpsBlue,draw=MumpsBlue] (6.25,3.75) circle (2.5pt);
    \node[anchor=south west] at (6.38,3.88)
      {$s_{pp}^{(k)}$};

    \filldraw[fill=MumpsCyan,draw=MumpsCyan] (8.35,3.75) circle (2.5pt);
    \node[anchor=north west] at (8.48,3.62)
      {$s_{pj}^{(k)}$};
      ;
  \end{tikzpicture}
  \caption{Bunch--Kaufman 选主元}
  \label{fig:bunch-kaufman-pivot-elements}
  \end{figure}

首先检查第 $k$ 列，记当前对角元大小为 $a_k=\left|s_{kk}^{(k)}\right|$，
其对角线以下最大元素为：
$$
\lambda_k=\max_{k<i\le f}\left|s_{ik}^{(k)}\right|
=\left|s_{pk}^{(k)}\right|,\qquad
p=\underset{k<i\le f}{\operatorname{argmax}}
  \left|s_{ik}^{(k)}\right|.$$
如果该列不全零，则按照以下规则选择主元：
\begin{enumerate}
  \item 若
  $a_k\ge\gamma\lambda_k$，
  直接采用 $s_{kk}^{(k)}$ 作为 $1\times1$ 主元，不进行交换。

  \item 若
  $a_k<\gamma\lambda_k$，
  则检查候选行 $p$。定义该行除对角元素之外的最大元素：
  $$\sigma_p=
  \max_{\substack{k\le j\le f\\j\ne p}}
  \left|s_{pj}^{(k)}\right|$$

  因为 $|s_{pk}^{(k)}|=\lambda_k$，所以 $\lambda_k>0$ 时有 $\sigma_p\ge\lambda_k>0$。随后依次判断：
  \begin{enumerate}[label=(\alph*)]
    \item 若
    $a_k\ge\gamma\frac{\lambda_k^2}{\sigma_p}$，
    仍然使用 $s_{kk}^{(k)}$ 作为 $1\times1$ 主元；

    \item 否则，若
    $\left|s_{pp}^{(k)}\right|\ge\gamma\sigma_p$，
    同时交换第 $k$、$p$ 行和第 $k$、$p$ 列，把 $s_{pp}^{(k)}$ 移到 $(k,k)$，并将其作为 $1\times1$ 主元；

    \item 若以上两个条件均不满足，选择由第 $k$、$p$ 个变量组成的 $2\times2$ 主元块。通过对称交换将变量 $p$ 移动到 $k+1$ 位置，得到：
    \[
    D_k=
    \begin{bmatrix}
    \widetilde s_{kk}^{(k)}&\widetilde s_{k,k+1}^{(k)}\\
    \widetilde s_{k+1,k}^{(k)}&\widetilde s_{k+1,k+1}^{(k)}
    \end{bmatrix}.
    \]
  \end{enumerate}
\end{enumerate}
对于 $1\times1$ 主元 $s_{pp}^{(k)}$，用置换矩阵 $P_k$ 同时交换第 $k$、$p$ 行和第 $k$、$p$ 列：
$$\widetilde S^{(k)}=P_k^TS^{(k)}P_k,
\qquad
D_k=\left[\widetilde s_{kk}^{(k)}\right].$$
对于由 $k$、$p$ 组成的 $2\times2$ 主元，用 $P_k$ 将这两个变量置换到活动子块的前两个位置：
$$\widetilde S^{(k)}=P_k^TS^{(k)}P_k
=
\begin{bmatrix}
D_k&B_k^T\\
B_k&C_k
\end{bmatrix},
\qquad
D_k\in\mathbb R^{2\times2}.$$
置换必须同时作用于整个波前矩阵的相应行和列，并同步更新全局变量编号。多步置换累积后，最终得到：
$$P=P_1P_2\cdots P_t,
\qquad
P^TFP=LDL^T.$$
完成对称置换后，
当前块列的消元乘子为：
$$L_k=B_kD_k^{-1}.$$
右下角活动子块通过对称 Schur 补更新：
$$S^{(k+r)}
=C_k-L_kD_kL_k^T
=C_k-B_kD_k^{-1}B_k^T.$$
对于 $1\times1$ 主元 $D_k=[d_k]$，公式退化为：
$$L_k=\frac{B_k}{d_k},
\qquad
S^{(k+1)}=C_k-\frac{B_kB_k^T}{d_k}.$$
对于 $2\times2$ 主元，更新为：
$$L_k=
\frac{1}{\Delta_k}
B_k
\begin{bmatrix}
c&-b\\
-b&a
\end{bmatrix},
\qquad
S^{(k+2)}=C_k-L_kD_kL_k^T.$$
这里就可以看出选择$2\times2$ 主元的精妙之处：\textbf{推导一下2x2的行列式}

但是仍然存在找不到列主元的情况：即当前列全零，则我们类似于非对称的处理将其移动到最后一个主元位置
，准备将其延迟到父节点的波前矩阵上再做处理。

\subsection{块分解优化}

与非对称情况相同，对称算法也使用分块方法优化。算法维护四个连续区间：

\begin{table}[htbp]
\centering
\caption{节点分区说明}
\label{tab:partition}
\small
\setlength{\tabcolsep}{8pt}
\renewcommand{\arraystretch}{1.16}
\begin{tabular}{@{}c >{\centering\arraybackslash}p{8.2cm}@{}}
\toprule
\textbf{区间} & \textbf{说明} \\
\midrule
$[0, \mathtt{panel\_begin})$
  & 以前面板已经消去 \\
$[\mathtt{panel\_begin}, k)$
  & 当前面板已经接受的主元 \\
$[k, \mathtt{active\_end})$
  & 仍可参与主元选择的活动节点 \\
$[\mathtt{active\_end}, n)$
  & 已经延迟的节点 \\
\bottomrule
\end{tabular}
\end{table}

设当前 panel 从
$b=\code{panel\_begin}$
开始。以前 panel 产生的 Schur 补已经全部写回 $A$，所以 panel 开始时
$A(b:n-1,b:n-1)$
就是最新的尾矩阵。

假设当前 panel 已经接受了 $t$ 列，令
$L_P=L(:,b:b+t-1)$，相应的 $1\times1$ 和 $2\times2$ 主元组成 $D_P$。定义工作矩阵
$W_P=L_PD_P$，当前真正的 Schur 补为 :
$$S=A_{\mathrm{stored}}-L_PD_PL_P^T
=A_{\mathrm{stored}}-L_PW_P^T.$$

准备处理第 $k$ 列时，$A(:,k)$ 尚未包含当前 panel 前面 $t$ 列主元的更新，因此必须先计算当前列的更新，再继续选择主元。

按照 Bunch--Kaufman 方法，首先判断当前列的对角元素是否满足阈值要求；
如果满足，则把它标记为主元。如果不满足，后续判定需要使用当前列最大元素所在行的数据，
因此还需要更新该行。
当 $1\times1$、$2\times2$ 主元均不满足要求时，需要延迟主元。
算法维护 \code{active\_end}，记录尾部延迟行、列的边界；每次延迟时，
从该边界向左寻找下一个非零列，再将其与被延迟主元交换。

每个 panel 的分解过程维护工作数组 $W_P$。理论上可以只保存 $L_P,D_P$，但每次需要当前列时，
都必须先计算
$D_PL_P(k,:)^T$，
再计算
$L_P\left(D_PL_P(k,:)^T\right)$。

由于 $D_P$ 混合了 $1\times1$ 和 $2\times2$ 块，反复执行这一过程较为复杂。预先保存
$W_P=L_PD_P$
后，当前列可直接计算
$L_PW_P(k,:)^T$，
只需要一次 \code{dgemv}，不必在每次更新时重新遍历 $D_P$ 的 $1\times1/2\times2$ 块结构。
\section{两种稠密矩阵分解算法的CUDA Kernel实现}